\documentclass[12pt,a4paper]{article}

\usepackage[T1]{fontenc}
\usepackage[utf8]{inputenc}
\usepackage[brazil]{babel}
\usepackage{graphicx}
\usepackage{geometry}
\usepackage{listings}
\usepackage{listingsutf8}
\usepackage{xcolor}
\usepackage{setspace}
\usepackage{float}
\usepackage{tikz}
\usepackage{tabularx}

\usetikzlibrary{positioning, shapes.geometric}
\onehalfspacing
\geometry{margin=2.5cm}

\lstset{
    language=SQL,
    basicstyle=\ttfamily\small,
    keywordstyle=\color{blue},
    commentstyle=\color{gray},
    stringstyle=\color{red},
    breaklines=true,
    showstringspaces=false
}

\newcommand{\montarinicio}[3]{
    % =========================
    % CAPA
    % =========================
    \begin{titlepage}
        \centering

        {\large Faculdade Anhanguera}\\[0.5cm]
        {\large Tecnólogo em DevOps}\\[3cm]

        {\Large \textbf{#1}}\\[2cm]

        {\large Gabriel Blank Stift Mousquer}\\[1cm]

        {\large Ivoti/RS}\\
        {\large 2026}

    \end{titlepage}

    % =========================
    % FOLHA DE ROSTO
    % =========================
    \begin{titlepage}
        \centering

        {\large Faculdade Anhanguera}\\[0.5cm]
        {\large Tecnólogo em DevOps}\\[2cm]

        {\Large \textbf{#1}}\\[1.5cm]

        \begin{flushright}
            Trabalho apresentado à disciplina de \\
            \textbf{#2} \\
            Professor: #3 \\
        \end{flushright}

        \vfill

        {\large Gabriel Blank Stift Mousquer}\\
        Matrícula: 2024201126\\

        \vfill

        {\large Ivoti/RS}\\
        {\large 2026}

    \end{titlepage}

    % =========================
    % SEÇÃO: TOC
    % =========================
    \tableofcontents
    \newpage
}


\begin{document}

\montarinicio
    {ALGORITMOS E PROGRAMAÇÃO ESTRUTURADA}
    {Algoritmos e Programação Estruturada}
    {Anderson Emídio De Macedo Gonçalves}

\section{Codificação completa}

O código-fonte desenvolvido em linguagem C é apresentado a seguir.

\lstset{language=C}

\lstinputlisting[
    caption={Código em C},
    inputencoding=utf8/latin1,
    basicstyle=\ttfamily\scriptsize
]{./static/main.c}

\newpage

\section{Explicação do algoritmo}

O programa tem como objetivo calcular o salário de um profissional que trabalha por hora, a partir do valor da hora trabalhada e da quantidade de horas trabalhadas no mês, considerando o cálculo do salário bruto, do desconto de 9\% e do salário líquido.

Inicialmente, o programa solicita ao usuário o valor da hora de trabalho (\texttt{valor\_hora}) e a quantidade de horas trabalhadas no mês (\texttt{qtde\_horas}). Esses dados são utilizados como entrada para os cálculos posteriores.

Em seguida, o cálculo do salário bruto é realizado pela função \texttt{calcular\_salario\_bruto()}, que recebe como parâmetros o valor da hora e a quantidade de horas trabalhadas. Essa função retorna o produto entre esses dois valores, resultando no salário bruto do profissional.

Com o salário bruto calculado, o programa chama a função \texttt{calcular\_desconto()}, que recebe o salário bruto como parâmetro e aplica uma taxa fixa de 9\%, definida pela constante \texttt{aliquota\_desconto}. O resultado dessa função corresponde ao valor do desconto aplicado sobre o salário.

Por fim, o salário líquido é calculado pela função \texttt{calcular\_salario\_liquido()}, que recebe o salário bruto e o valor do desconto como parâmetros. Essa função realiza a subtração entre o salário bruto e o desconto, resultando no valor final a ser recebido pelo profissional.

Após a realização dos cálculos, o programa exibe na tela o salário bruto, o valor do desconto e o salário líquido, todos formatados com duas casas decimais.

A implementação garante que:
\begin{itemize}
    \item O programa leia corretamente o valor da hora e a quantidade de horas trabalhadas;
    \item O salário bruto seja calculado corretamente a partir da multiplicação dos valores de entrada;
    \item O desconto de 9\% seja aplicado corretamente sobre o salário bruto;
    \item O salário líquido seja obtido corretamente pela subtração do desconto;
    \item Os três valores finais sejam exibidos de forma clara ao usuário;
    \item O programa utilize funções separadas para cada etapa do cálculo, garantindo modularização.
\end{itemize}

\newpage

\section{Execução do código}

Foram realizados testes com diferentes conjuntos de valores para validar o funcionamento do programa.

\subsection{Caso 1: Valores redondos}

\begin{figure}[H]
    \centering
    \includegraphics[width=0.8\textwidth]{./static/img-1.png}
\end{figure}

\subsection{Caso 1: Valor da hora decimal}

\begin{figure}[H]
    \centering
    \includegraphics[width=0.8\textwidth]{./static/img-2.png}
\end{figure}

\subsection{Caso 1: Quantidade de horas decimal}

\begin{figure}[H]
    \centering
    \includegraphics[width=0.8\textwidth]{./static/img-3.png}
\end{figure}

\end{document}
